\section{Discussion and policy implications}
\label{sec:discussion}

The results should be read as scenario analysis, not prediction. We do not forecast the pace or extent of AI adoption in the United Kingdom; we trace the distributional and fiscal consequences that would follow \emph{if} labour-market shocks of a given size and occupational incidence were to materialise, under tax-benefit rules as legislated. The early observational record justifies the full width of the grid, including its benign corner: population-scale Danish evidence finds minimal earnings and hours effects of chatbot adoption to date \citep{humlum2025}, and postings-based analyses disagree on whether the entry-level slowdown in exposed occupations is attributable to AI at all. UK postings evidence points to a modest relative contraction in exposed occupations since late 2022 \citep{henseke2025}, but nothing yet at the scale of our central scenario. The Office for Budget Responsibility treats AI principally as a productivity risk to its fiscal projections \citep{obr2025}; our estimates supply the complementary labour-incidence arithmetic that such projections currently leave implicit. The scenario parameters are calibration conventions transported from headline estimates in the empirical and modelling literature---including \citet{briggs2023} and \citet{cazzaniga2024}---rather than elasticities estimated on UK data, and the wide span between our mildest and harshest cells is a deliberate acknowledgement of the uncertainty surrounding all of these anchors. The incidence gradient itself is an assumption in every family we run; what the paper establishes is how the tax-benefit system mediates each assumed incidence.

Two findings organise the discussion. The first is that the UK tax-benefit system provides substantial automatic insurance to low-income households: means-tested support---principally Universal Credit---absorbs a large share of the market-income losses in the lower deciles, while progressive income taxation claws back part of the losses higher up before they reach disposable income. This cushioning is not free. In the central scenario the combined cost of lower receipts and higher benefit expenditure is \pounds 18.4~billion a year, and the cost scales with the severity of the shock. The protection households experience is purchased through a deterioration of the public finances.

The second finding is that the Gini of equivalised household disposable income rises in every one of the 66 shock-size cells and in all five incidence families (Sections~\ref{sec:results-grid} and~\ref{sec:results-incidence}), even though the market-income incidence in the exposure-proportional family is progressive---displacement risk rises monotonically with household income (Figure~\ref{fig:transition}), and average losses are largest at the top. That tension is the interesting result. Each displaced worker falls from somewhere on the earnings ladder to the bottom of it, while wage and capital gains lift those who remain; because the displaced come disproportionately from higher-income households, the shock simultaneously thins the upper-middle of the distribution and swells its lower tail. When households are re-ranked on post-shock income, as the Gini requires, dispersion increases even though decile-average losses computed at fixed baseline ranks look progressive. A similar re-ranking mechanism is documented for Ireland by \citet{doorley2026}, and it distinguishes the AI case from earlier automation waves, in which exposure fell on mid- and low-wage routine work \citep{autor2003, doorley2023automation}.

\subsection*{Policy implications: what depends on who bears the shock}

The central policy message follows from the incidence axis rather than from any single scenario. Inequality rises whichever family one assumes; what the incidence determines is whether the UK faces primarily a fiscal-resilience problem or a poverty-protection problem.

If the shock is top-loaded, as in the exposure-proportional family, the problem is chiefly fiscal. Displaced workers are drawn from occupations that pay well and are taxed heavily, so each displacement removes a large slug of income tax and National Insurance while adding comparatively little to means-tested expenditure. The Exchequer cost of the central shock is \pounds 18.4~billion under exposure-proportional incidence against \pounds 14.2~billion under uniform incidence---the same aggregate shock, roughly 30 per cent more expensive purely because of who bears it. Under top-loaded incidence the first-order policy question is revenue resilience: the UK tax base leans heavily on the taxation of labour income, and the workers most exposed are precisely those who contribute most to it. If AI additionally shifts the functional distribution of income from labour towards capital, the base that finances the cushioning we document would itself erode---a possibility that has motivated recent work on the taxation of capital in an AI-intensive economy \citep{bastani2024}. The tax-composition exercise in Section~\ref{sec:policy} turns that possibility into a quantified scenario for the Treasury: because the effective labour-tax rate on displaced earnings is almost exactly the corporation-tax main rate, the net revenue shortfall is governed by the share of the lost wage bill that reappears as taxable corporate profit---full reappearance is roughly revenue-neutral, no reappearance leaves the whole labour-tax loss uncovered, and the harm is distributional rather than fiscal, since recycled dividends raise the Gini while reaching no household near the poverty line. The exercise is static accounting on top of the microsimulation, with corporation-tax incidence not modelled, but it gives base-erosion discussions a concrete UK arithmetic to anchor on.

If instead the shock lands flatter or on junior workers, the fiscal exposure moderates but the poverty exposure does not. Uniform incidence is the cheapest family we run, yet its poverty effect (+1.83 points before housing costs) is essentially the same as the exposure family's (+1.87), and its Gini effect is the largest of the five families (+1.14 points). Junior-concentrated incidence costs slightly more than the exposure family (\pounds 19.4~billion) with a near-identical poverty effect. In these worlds the binding constraint is the adequacy of out-of-work support, and the relevant instruments are benefit levels, the contributory window, and---for the junior case---the entry prospects of young workers. Evidence that generative AI operates as a seniority-biased technology, compressing demand for junior and entry-level roles in exposed occupations \citep{hosseini2026, brynjolfsson2025}, with early UK corroboration in \citet{kleinteeselink2025}, argues for close monitoring of youth labour-market entry. Our junior-concentrated family can only re-weight who is displaced within a single year; the dynamic costs of a scarred entry cohort---foregone experience, delayed progression, persistent earnings losses---lie outside a static framework, so even that family understates the long-run burden on new entrants.

The fiscal-versus-poverty framing also now has a policy answer rather than only a diagnosis. The counterfactual reforms in Section~\ref{sec:policy} show that temporary Universal-Credit-side stabilisers---a circuit-breaker uplift to the standard allowance, or a benefit-cap suspension combined with a taper cut---deliver a point of after-housing-costs poverty reduction at roughly \pounds 10~billion, against \pounds 13--14~billion for an earnings-linked wage-insurance scheme, because wage insurance is earnings-proportional and much of its benefit accrues in the upper half of the baseline distribution. Wage insurance buys more total poverty reduction, but at disproportionate cost; a government facing a top-loaded shock and a strained revenue base gets more protection per pound by working through the existing means-tested architecture than by building an earnings-replacement scheme alongside it.

The expertise-compression family, our stylised stress test of the view that AI erodes the returns to accumulated expertise, is the worst of both worlds among the stylised families: the most expensive (\pounds 24.5~billion) and the most poverty-raising (+2.12 points) of the four---and the measured family exceeds it on both margins (Table~\ref{tab:incidence}). If that view proves correct, neither fiscal buffers nor poverty protection alone suffices; both margins bind at once.

The shock also has a geography, and it does not track existing regional-policy maps. In the constituency analysis of Section~\ref{sec:results-geo}, the largest proportionate income losses fall on Northern Ireland---the worst-hit region, at $-2.3$ per cent of aggregate household income---and Scotland, while Wales and the South East are mildest; the hardest-hit constituencies mix Northern Irish seats with inner-London ones. An AI labour shock is not a rerun of deindustrialisation: because exposure follows clerical and professional employment, it strikes administrative and professional centres wherever they sit, and a levelling-up strategy calibrated to the geography of manufacturing decline would miss much of it. These geographic results rest on imputed occupations on the enhanced 2023--24 dataset (Appendix~\ref{app:geo-imputation}) and are indicative of relative, not absolute, local impacts.

Two implications hold across all families. First, because the modelled losses fall on identifiable occupational groups rather than diffusely across the workforce, retraining and lifelong-learning provision targeted at exposed occupations is the most direct lever available; the tax-benefit system can replace income but cannot restore occupational attachment. Second, since measured inequality rises under every incidence we examine, distributional monitoring should not wait for evidence on which incidence is materialising: the direction of the inequality effect is not in doubt within this framework, only its composition.

Taken together, the results suggest that the United Kingdom would enter a period of AI-driven labour-market disruption with a tax-benefit system that redistributes the losses effectively but does not, and cannot, prevent them. The system converts each market-income shock into a smaller disposable-income shock at a substantial, incidence-dependent fiscal cost, while the polarising character of displacement pushes measured inequality upwards regardless of who bears it. Policy attention is best directed at the margins the system does not reach: the occupational reallocation of exposed workers, the entry prospects of the young, and the robustness of a labour-taxation-heavy revenue base.

\section{Limitations and further work}
\label{sec:limitations}

Several limitations qualify the results. We state them plainly, together with the measurement conventions a reader needs in order to compare our figures with official statistics.

\paragraph{Measurement conventions.} All income results are on the HBAI cash disposable income concept at the household level, equivalised; poverty and inequality statistics share that concept. In computing the Gini we bottom-code negative equivalised incomes at zero. Our baseline Gini is 0.309, below the official UK figure of roughly 0.34, because we use the plain Family Resources Survey without the SPI adjustment that corrects for under-coverage of top incomes; levels should therefore not be compared directly with official statistics, though changes---the objects of interest---are computed on a consistent basis throughout.

\paragraph{Monte Carlo conventions.} Which results rest on which random draws matters for interpretation; the conventions are set out in Section~\ref{sec:microsim} and Appendix~\ref{app:seeds}. Single-draw figures carry draw noise that the averaged figures do not, although the 20-draw runs indicate the aggregate magnitudes are stable across seeds.

\paragraph{Exposure resolution and crosswalk provenance.} Occupational exposure is assigned at the 1-digit SOC level, because the Family Resources Survey records occupation only at the major-group level; within-group heterogeneity in AI exposure---substantial in the indices of \citet{felten2021}, \citet{webb2019} and \citet{eloundou2023}---is averaged away. The exposure scores themselves reach the FRS through a reconstructed pipeline: the complementarity adjustment is rebuilt from published sources (recovering the published ranking up to a level offset), and the US occupational scores are carried to SOC2020 through a chained SOC2018--SOC2010--ISCO--SOC2020 crosswalk with ASHE employment weights available for 340 of 412 unit groups. A UK-native task-based index (GAISI) built from British Skills and Employment Survey data now exists \citep{henseke2025}; adopting it, together with an imputation of detailed occupation from the Labour Force Survey, would remove the crosswalk and the resolution constraint at once, and is the first item on our agenda.

\paragraph{Survey weights.} We use the uncalibrated FRS grossing weights rather than PolicyEngine's enhanced dataset. The enhanced dataset improves aggregate fiscal accuracy through reweighting and household cloning, but cloning breaks the one-to-one correspondence between records and observed occupations on which the exposure join depends. Aggregate fiscal magnitudes should be read with the usual caveats attaching to unadjusted survey grossing.

\paragraph{Static displacement.} Displaced workers are fully out of work in the simulated year---earnings, hours, employee pension contributions and salary sacrifice set to zero, no statutory pay, benefit entitlements recomputed under existing rules---while dual jobholders retain self-employment income. The duration structure of unemployment is not modelled, so the exhaustion of contributory benefits (new-style Jobseeker's Allowance is payable for at most six months) cannot be expressed; over longer horizons we overstate the support available to displaced workers without means-tested entitlements.

\paragraph{The capital channel.} The capital-income measure covers only interest and dividends, excluding rental income and returns inside private pensions, and household surveys under-represent top incomes and capital income in particular. The concentration of capital income in the top decile that drives this channel is, if anything, understated. Linked administrative or wealth-survey data (the Wealth and Assets Survey) would sharpen it considerably.

\paragraph{The self-employed.} All shock scenarios exclude the self-employed. Exposure indices and the displacement literature underpinning our parameters are defined over employees, and self-employed earnings dynamics under AI adoption may include both displacement and new opportunities; we exclude the group rather than impose an arbitrary treatment. Since self-employment is a meaningful share of UK employment and is concentrated in some exposed occupations, this is a substantive omission.

\paragraph{Mechanisms and anchors.} Where we attribute decile-level patterns to particular mechanisms---the thinness of savings buffers at the bottom, income pooling within couples---these attributions are hypotheses consistent with the data rather than results of a formal decomposition, and we flag them as such. The scenario anchors, likewise, are author-designed stress tests transported from headline estimates: the 1 per cent ``low'' cell is a sensitivity case rather than an employment estimate, the 7 per cent displacement and 2.6 per cent wage figures are conventions converting task-exposure and productivity estimates into single-year shocks, and the paper-anchored scenarios are transpositions, not implementations, of the studies they draw on. As UK-specific evidence accumulates---early examples include \citet{kleinteeselink2025} and \citet{rockall2025}---the scenarios can be re-anchored, and the gap between scenario analysis and conditional forecast will narrow. Extending the exposure measurement, restoring calibrated weights through an occupation-preserving reweighting, and introducing benefit-duration dynamics are the immediate priorities for further work.
